\chapter{The proper Riemann integral}
\SourceRef{Zorich I, Section 6.1, PDF pp. 351--363, for mesh-tagged sums; Rudin, Chapter 6, for Darboux sums with integrator $x$; Apostol, Section 7.26, PDF pp. 182--185, for countable covers and oscillation; Tao I, Chapter 11, for step-function bounds}
The basic real definition is bounded Darboux integration on a compact
nondegenerate interval. Integrability is a separate predicate, not a consequence
of writing an integral symbol or assigning a default value to nonintegrable data.

\KEntry{K03.08.partitions}{Partitions and tags}
\begin{definition}
For $a<b$, a partition is $P=(a=x_0<x_1<\cdots<x_n=b)$.
Its mesh is $|P|=\max_i\Delta x_i$, where $\Delta x_i=x_i-x_{i-1}$.
A tagging chooses $\xi_i\in[x_{i-1},x_i]$, and its sum is
$S(f;P,\xi)=\sum_i f(\xi_i)\Delta x_i$. A refinement inserts partition
points. Two partitions have the common refinement formed by their union.
For $[a,a]$ we use the empty sum, equal to zero, and no positive mesh.
\end{definition}

\KEntry{K03.08.darboux_sums}{Upper and lower sums}
\begin{definition}
For bounded $f:[a,b]\to\R$, set
\[M_i=\sup_{[x_{i-1},x_i]}f,\quad m_i=\inf_{[x_{i-1},x_i]}f,\quad
 U(f,P)=\sum_iM_i\Delta x_i,\quad L(f,P)=\sum_im_i\Delta x_i.
\]
\end{definition}
\begin{lemma}
Refinement increases $L$ and decreases $U$. For arbitrary $P,Q$,
$L(f,P)\le U(f,Q)$, and every tagged sum for $P$ lies between its
lower and upper sums.
\end{lemma}
\begin{proof}
On a subinterval inserted into an old cell, its infimum is at least the old
infimum and its supremum at most the old supremum. Multiplication by the
positive new lengths and their additivity gives the refinement inequalities.
For the common refinement $R$, $L(P)\le L(R)\le U(R)\le U(Q)$.
The last assertion is the termwise inequality $m_i\le f(\xi_i)\le M_i$.
\end{proof}

\KEntry{K03.08.integral}{The integrability predicate}
\begin{definition}
The lower and upper integrals are
\[\underline{I}(f)=\sup_P L(f,P),\qquad
\overline{I}(f)=\inf_P U(f,P).
\]
Boundedness ensures that these are finite and $\underline I\le\overline I$.
The function is Riemann integrable if they are equal; their common value
is $\int_a^bf$. A complex function is integrable if its real and imaginary
parts are, and its integral is their componentwise integral.
\end{definition}
For comparison with step-function definitions, a partition provides step
minorants and majorants with constants $m_i,M_i$ on the open cells,
assigning the finitely many endpoints their actual values. Their integrals
are $L(P),U(P)$. Conversely, for a finite step minorant $s\le f$, refine
at its finitely many breaks. Off those breaks its constants are bounded by
the infima of $f$ on the open cells. Shrinking the cells adjacent to breaks
changes a sum by at most a fixed bound times their total length. Consequently
$\int s\le\underline I(f)$ on letting that length tend to zero. The same
argument for majorants gives $\overline I(f)\le\int t$ for $t\ge f$.
Taking suprema and infima, together with the partition-provided bounds,
proves equality of the two definitions. Endpoint values have zero contribution
because finitely many neighboring lengths can be made arbitrarily small;
this is a proof of the comparison, not an identification of notation.

\KEntry{K03.08.criterion}{The Darboux criterion}
\begin{theorem}
A bounded real function is integrable if and only if for each $\eps>0$
some partition $P$ satisfies $U(f,P)-L(f,P)<\eps$.
\end{theorem}
\begin{proof}
If $\underline I=\overline I=I$, choose $P,Q$ with
$U(P)<I+\eps/2$ and $L(Q)>I-\eps/2$. Their common refinement
has gap less than $\eps$. Conversely every such partition gives
$0\le\overline I-\underline I\le U(P)-L(P)<\eps$, hence equality.
\end{proof}

\KEntry{K03.08.tagged_equivalence}{Equivalence with all fine tagged sums}
\begin{theorem}
For bounded real $f$, Darboux integrability with value $I$ is equivalent to:
for every $\eps>0$ there is $\delta>0$ such that
$|S(f;Q,\xi)-I|<\eps$ for every tagging of every $Q$ with $|Q|<\delta$.
The analogous statement for complex functions is equivalent componentwise.
\end{theorem}
\begin{proof}
Let $|f|\le M$ and fix $P$ with Darboux gap $<\eps/2$.
If $P$ has $m$ interior points, cells of $Q$ whose interiors contain
such a point have total length at most $m|Q|$; even counting cells
touching such points gives the harmless bound $2m|Q|$.
Every other $Q$ cell lies in a single $P$ cell. Insert the $P$ points
into $Q$ and compare sums: the good terms are unchanged, while the total
change on the exceptional cells has magnitude at most $4Mm|Q|$.
The refined tagged sum is between $L(P)$ and $U(P)$, as is $I$.
Thus $|S(Q,\xi)-I|\le\eps/2+4Mm|Q|$. Choose $\delta$ to make the
second term $<\eps/2$; if $M m=0$ it is already zero.

Conversely suppose all fine sums approach $I$. Choose a fine $Q$ such that
every sum differs from $I$ by less than $\eps/4$. In each cell choose a
tag with value within $\eps/(8(b-a))$ of its supremum and, independently,
a tag equally close to its infimum. Such tags exist by the definition of
supremum, without assuming extrema are attained. The two sums differ by
at most $\eps/2$, and together approximate $U(Q),L(Q)$ with total error
at most $\eps/4$. Hence $U(Q)-L(Q)<\eps$. The criterion proves
integrability, and the first direction plus uniqueness of a limit proves
the value is $I$. For complex values use $|\Re z|,|\Im z|\le|z|$ in one
direction and $|z|\le|\Re z|+|\Im z|$ in the other.
\end{proof}

\KEntry{K03.08.cauchy}{Cauchy criterion for sums}
\begin{theorem}
For bounded scalar $f$, a common fine-mesh limit exists if and only if
for every $\eps>0$ all two tagged sums of mesh below some $\delta$
differ in norm by less than $\eps$.
\end{theorem}
\begin{proof}
A common limit implies the assertion by the triangle inequality.
Conversely choose any sequence of tagged partitions with meshes tending
to zero (uniform partitions suffice). Their sums are a Cauchy sequence,
so scalar completeness gives a limit $I$. Given $\eps$, choose the
threshold for mutual distance $\eps/2$, then one of these sums below that
threshold and within $\eps/2$ of $I$. Every sum of mesh below the same
threshold is within $\eps$ of $I$. The preceding equivalence identifies
this limit with the Darboux integral.
\end{proof}

\KEntry{K03.08.continuous}{Continuous scalar functions}
\begin{theorem}
Every continuous real or complex function on $[a,b]$ is integrable.
\end{theorem}
\begin{proof}
Compactness gives boundedness and uniform continuity. For real $f$, choose
$\delta$ so that $|f(x)-f(y)|<\eps/(b-a)$ when $|x-y|<\delta$.
Any partition of mesh below $\delta$ then has each oscillation
$M_i-m_i\le\eps/(b-a)$, and hence gap at most $\eps$.
Apply the criterion with a smaller input tolerance for a strict bound.
For complex $f$, its two continuous real components are integrable.
\end{proof}

\KEntry{K03.08.monotone}{Monotone, step and piecewise continuous functions}
\begin{theorem}
Bounded monotone real functions, finite step functions, and functions
continuous on each of finitely many open cells with finite one-sided
boundary limits are integrable, provided all assigned boundary values are finite.
\end{theorem}
\begin{proof}
For increasing $f$ and a uniform partition of mesh $(b-a)/n$,
\[U-L=\frac{b-a}{n}\sum_{i=1}^n(f(x_i)-f(x_{i-1}))
=\frac{b-a}{n}(f(b)-f(a))\longrightarrow0.
\]
Decreasing functions follow on changing sign. A step function has zero
oscillation outside cells touching its finitely many breaks; their total
length is at most twice the number of breaks times the mesh, making the
gap tend to zero. A function in the last class extends continuously to
each closed cell using its one-sided limits. Those finitely many continuous
pieces are integrable by the previous theorem, and changing finitely many
endpoint values has zero effect by \KRef{K03.08.finite_changes}.
Combining their partitions proves integrability on the whole interval;
the same partition argument proves the additivity used here.
\end{proof}

\KEntry{K03.08.linearity}{Linearity and splitting intervals}
\begin{theorem}
Finite scalar linear combinations of integrable functions are integrable
and their integrals are the same combinations. For $c\in[a,b]$, a bounded
function is integrable on $[a,b]$ if and only if both restrictions to
$[a,c]$ and $[c,b]$ are, and then
$\int_a^bf=\int_a^cf+\int_c^bf$.
\end{theorem}
\begin{proof}
For real functions the oscillation inequality
$\osc(\alpha f+\beta g,J)\le|\alpha|\osc(f,J)+|\beta|\osc(g,J)$,
on common refinements of their small-gap partitions, proves integrability.
Every tagged sum is linear, so its limit is linear. For splitting, insert
$c$ into a partition. Its Darboux gap is the sum of the two nonnegative
restricted gaps. Small whole gaps therefore give small restricted gaps;
conversely combine two small-gap partitions. Tagged sums for partitions
containing $c$ split exactly, and their limits give the asserted equality.
Complex versions follow by components and complex scalar multiplication.
\end{proof}

\KEntry{K03.08.order}{Order, modulus and norm estimates}
\begin{theorem}
For integrable real $f\le g$, $\int f\le\int g$. For integrable real
or complex $f$, $|f|$ is integrable and
\[\left|\int_a^bf\right|\le\int_a^b|f|\le(b-a)\sup_{[a,b]}|f|.
\]
\end{theorem}
\begin{proof}
All paired tagged sums preserve a real pointwise inequality, so their
limits preserve it. For real values $||u|-|v||\le|u-v|$ bounds the
Darboux gap for $|f|$ by that for $f$. For complex values the oscillation
of $|f|$ is at most the sum of those of the two components, so common
partitions prove its integrability as well. Put $z=\int f$. If $z\ne0$,
choose $\lambda=\bar z/|z|$. Then
$|z|=\int\Re(\lambda f)\le\int|f|$; if $z=0$ the assertion follows
from positivity. Comparing $|f|$ with its supremum proves the last bound.
\end{proof}

\KEntry{K03.08.products}{Products and continuous composition}
\begin{theorem}
Products of bounded integrable scalar functions are integrable. If $f$ is
real integrable and $\Phi$ is continuous on a compact interval containing
its range, then $\Phi\circ f$ is integrable. In particular $1/f$ is
integrable when $|f|\ge d>0$.
\end{theorem}
\begin{proof}
If $|f|\le A$, $|g|\le B$, then
$\osc(fg,J)\le A\osc(g,J)+B\osc(f,J)$; take common small-gap
partitions. Complex products reduce to four real products.
For composition let $|\Phi|\le H$. Given $\eta>0$, uniform continuity
provides $d>0$ such that $|u-v|<d$ gives
$|\Phi(u)-\Phi(v)|<\eta$. On cells with $\osc(f,J)<d$ the new
oscillation is at most $\eta$. The total length of the remaining cells
is at most $(U(f,P)-L(f,P))/d$. Thus the composition gap is at most
$\eta(b-a)+2H(U-L)/d$, which can be made arbitrarily small in that order.
For reciprocals, even when the range crosses both signs, use
$|1/u-1/v|\le|u-v|/d^2$ directly on the actual range. The same estimate
works for a complex reciprocal with modulus at least $d$.
\end{proof}

\KEntry{K03.08.finite_changes}{Finite changes}
\begin{theorem}
If bounded $f,g$ differ at only finitely many points and $f$ is integrable,
then $g$ is integrable with the same integral.
\end{theorem}
\begin{proof}
It suffices to treat $h=g-f$, supported on $N$ points, with $|h|\le M$.
Only partition cells touching those points contribute nonzero oscillation;
their total length is at most $2N|P|$. Thus the Darboux gap is at most
$4MN|P|\to0$, and every tagged sum has magnitude at most
$2MN|P|\to0$. Hence $h$ is integrable with value zero, proving the
claim by linearity, componentwise if necessary. This argument also
justifies all finite endpoint assignments in step-function comparisons.
\end{proof}

\KEntry{K03.08.uniform_stability}{Uniform limits preserve integrability}
\begin{theorem}
If integrable $f_n:[a,b]\to\K$ converge uniformly to $f$, then $f$ is
bounded and integrable, and
$\int f_n\to\int f$ with error at most $(b-a)\|f_n-f\|_\infty$.
\end{theorem}
\begin{proof}
A uniform error at most one from any bounded $f_N$ bounds $f$.
In the real case,
$U(f,P)-L(f,P)\le U(f_N,P)-L(f_N,P)+2(b-a)\|f-f_N\|_\infty$.
First choose $N$ with small uniform error and then a small-gap partition
for $f_N$. This proves integrability; the norm estimate of the previous
entry gives convergence of integrals. Apply the argument to components
for complex values, and then use the complex integral norm estimate for
the same final error bound.
\end{proof}

\KEntry{K03.08.zero_covers}{Length-zero sets using countable covers}
\begin{definition}
$E\subset\R$ has length zero if, for each $\eps>0$, it is covered by
countably many open bounded intervals $I_j$ with $\sum_j|I_j|<\eps$.
Subsets of length-zero sets have length zero.
\end{definition}
\begin{theorem}
Countable unions of length-zero sets have length zero. No nondegenerate
interval has length zero.
\end{theorem}
\begin{proof}
For $E=\bigcup_{n\ge1}E_n$, cover $E_n$ with intervals of total length
less than $\eps2^{-n-1}$. Enumerate the pairs of indices in the resulting
countable family. Finite sums of lengths are bounded by the sum of the
row bounds, so the total is less than $\eps$.
If a countable open cover covers $[c,d]$, compactness supplies a finite
subcover. Its total length is at least $d-c$: sort all its endpoints
together with $c,d$; each resulting cell inside $[c,d]$ is covered by
at least one interval, and summing its length counts at most the total
length of the covering intervals. Thus $[c,d]$ cannot have length zero
when $c<d$. Every nondegenerate interval contains such a compact interval.
\end{proof}

\KEntry{K03.08.zero_vs_content}{Countable length zero versus finite content}
\begin{example}
Enumerate $\Q\cap[0,1]=\{r_n:n\ge1\}$. Intervals centered at $r_n$
of length $\eps2^{-n-1}$ cover this set with total length less than
$\eps$, so it has length zero. However, any finite interval cover of it
has total length at least one. Indeed the union of the closures of those
finitely many intervals is closed and contains a dense subset of $[0,1]$,
hence contains all of $[0,1]$. The elementary finite-length inequality
just proved gives the lower bound one. Thus countable-cover length zero
cannot be replaced by arbitrarily small finite-cover content.
\end{example}

\KEntry{K03.08.cantor_set}{The middle-thirds set}
\begin{construction}
Let $C_0=[0,1]$, and obtain $C_{n+1}$ by deleting the open middle third
of each component of $C_n$. Put $C=\bigcap_{n\ge0}C_n$.
\end{construction}
\begin{proof}[Compactness, perfection, empty interior and length zero]
$C_n$ consists of $2^n$ disjoint closed intervals of length $3^{-n}$.
They are nested, so $C$ is closed in $[0,1]$, compact and nonempty
(it contains $0,1$). Every component endpoint remains in all later sets.
For $x\in C$ and every $n$, the component of $C_n$ containing $x$ has
two endpoints in $C$, at least one different from $x$, at distance at most
$3^{-n}$. Hence $x$ is a limit point of $C$: the set is perfect.
If an open interval were contained in $C$, it would be contained in every
$C_n$, impossible once its length exceeds $3^{-n}$, since a connected
interval cannot cross the gaps between distinct components. Thus the
interior is empty. The total length of the $C_n$ components is $(2/3)^n$.
Enlarge their finitely many closed intervals slightly to open intervals,
adding less than any prescribed further length. Choosing $n$ large
produces arbitrarily small open covers, so $C$ has length zero.
Equivalently, $C$ consists of numbers admitting a ternary expansion using
only digits $0,2$: membership chooses a left or right component at each
stage, whose nested intervals shrink to that digit expansion, and conversely
such digits always select surviving components. At a deleted endpoint
use the expansion with only $0,2$, for example $1/3=0.0222\ldots_3$;
the alternative $0.1000\ldots_3$ is not the membership convention.
\end{proof}

\KEntry{K03.08.oscillation_criterion}{The discontinuity criterion}
\begin{theorem}[Lebesgue criterion, proved with interval covers]
A bounded real function on $[a,b]$ is Riemann integrable if and only if
its set of discontinuities has length zero.
\end{theorem}
\begin{proof}
For $x\in[a,b]$ define its relative oscillation by
\[\omega(x)=\inf_{r>0}\osc(f,(x-r,x+r)\cap[a,b]).\]
It is zero exactly when $f$ is continuous at $x$: continuity makes all
nearby values close, and conversely oscillation below $\eps$ bounds
$|f(y)-f(x)|$ for all sufficiently close $y$. The set
$E_\alpha=\{x:\omega(x)\ge\alpha\}$ is closed, since a neighborhood
with oscillation less than $\alpha$ also supplies such a neighborhood
at every sufficiently close point. It is consequently compact.
The discontinuity set is $D=\bigcup_{n\ge1}E_{1/n}$.

Suppose $f$ integrable. Choose a partition with gap below $\alpha\eta$.
Every $x\in E_\alpha$ in the interior of a cell forces that cell's
oscillation at least $\alpha$. The total length of these bad cells
is therefore below $\eta$. Their interiors cover $E_\alpha$ except
possibly for finitely many partition endpoints. Enlarge the bad cells
and cover those endpoints by tiny open intervals, adding less than $\eta$
to total length. Since $\eta$ is arbitrary, $E_\alpha$ has length zero,
and the countable-union theorem makes $D$ length zero.

Conversely suppose $D$ length zero and $|f|\le M$.
Choose $\alpha>0$ and cover the compact $E_\alpha\subset D$ by finitely
many open intervals of total length $<\eta$; let their union be $G$.
For every $x\notin E_\alpha$ select a relative open neighborhood $V_x$
on which oscillation is $<\alpha$. Together with $G$ these cover $[a,b]$.
There is a mesh $\delta>0$ such that every cell of mesh below $\delta$
is contained in one member of this cover. To see this elementary cover
fact, choose about each point a ball whose doubled radius lies in one
cover member, take finitely many of the smaller balls covering the compact
interval, and use the least of their radii. A cell of smaller length
meeting a selected smaller ball lies in its doubled ball.
On cells contained in a $V_x$ the oscillation is $<\alpha$.
All other cells lie in $G$, and their total length is at most $\eta$:
their disjoint interiors form a subset of the finite union of covering
intervals, so the finite endpoint length calculation applies.
Hence $U-L\le\alpha(b-a)+2M\eta$. Choose $\alpha$ and then $\eta$
arbitrarily small. The Darboux criterion proves integrability.
This argument constructed no Lebesgue integral and used no measure-theoretic
convergence theorem.
\end{proof}

\KEntry{K03.08.dirichlet}{The rational indicator}
\begin{example}
Let $D(x)=1$ on rational $x\in[0,1]$ and $D(x)=0$ otherwise.
Every nondegenerate interval contains both a rational and an irrational:
Archimedean rational density gives the first, and translating a sufficiently
small nonzero irrational multiple into a rational subinterval gives the
second. Thus every Darboux cell has supremum one and infimum zero.
For every partition $U(D,P)=1$, $L(D,P)=0$, so $D$ is not integrable.
It is discontinuous at every point for the same density reason.
\end{example}

\KEntry{K03.08.thomae}{Thomae's function}
\begin{example}
For a reduced fraction $p/q\in[0,1]$ with $q>0$ set $T(p/q)=1/q$,
and put $T(x)=0$ on irrationals. It is continuous at every irrational,
discontinuous at every rational, and integrable with integral zero.
\end{example}
\begin{proof}
For fixed $N$, only finitely many reduced fractions in $[0,1]$ have
$q\le N$, since $0\le p\le q$. Near an irrational $x$, choose a
neighborhood avoiding this finite set; then $0\le T<1/N$ there, proving
continuity. At a rational $p/q$, irrational points approaching it have
value zero, different from $1/q$, so continuity fails (one-sided at
endpoints). For an integral proof independent of the discontinuity criterion,
put tiny intervals around the finitely many rationals of denominator
$\le N$ and choose a partition for which all cells touching them have
total length less than $\eta$. On every other cell $T\le1/N$.
Thus $U(T,P)\le\eta+1/N$. Every cell contains an irrational, so
$L(T,P)=0$. Let $N\to\infty$ and $\eta\downarrow0$ to obtain the result.
\end{proof}

\KEntry{K03.08.countable_changes}{Countable modifications can fail}
\begin{example}
The zero function is integrable with value zero, while the rational
indicator is not integrable. They differ exactly on the countable set
$\Q\cap[0,1]$, which has length zero. Thus equality off a length-zero
set does not by itself preserve Riemann integrability. If both functions
are already bounded and Riemann integrable, such equality does imply
equal integrals: their difference $h$ is zero off a length-zero set;
at every continuity point of $h$, density of that complement forces
$h=0$. Alternatively, cover its level sets by the oscillation argument.
For a direct value proof, every cell has a point outside the zero set,
so one can choose tags there to make all sums for $h$ equal zero; the
all-tagged limit then equals zero. The distinction is the separate
existence hypothesis on both integrals.
\end{example}

\KEntry{K03.08.oriented}{Orientation and finite-dimensional values}
\begin{definition}
For $a>b$ define $\int_a^bf=-\int_b^af$, requiring integrability on
$[b,a]$; define $\int_a^af=0$. For a finite-vector function
$f=(f_1,\ldots,f_d)$, integrability means integrability of every component,
and $\int f=(\int f_1,\ldots,\int f_d)$.
\end{definition}
\begin{proposition}
Interval additivity holds for arbitrary ordered triples whenever $f$ is
integrable on the compact interval containing them. For Euclidean norm,
$\|\int_a^bf\|\le\int_{\min(a,b)}^{\max(a,b)}\|f\|$.
\end{proposition}
\begin{proof}
Sort the three points; ordinary additivity on the resulting adjacent
intervals and the sign convention prove every ordering. The norm of $f$
is integrable because its oscillation is bounded by the sum of component
oscillations. If $z=\int f\ne0$, the real linear functional
$v\mapsto\langle z/\|z\|,v\rangle$ has norm one and takes value
$\|z\|$ at $z$. Componentwise linearity and the pointwise bound by
$\|f\|$ prove the estimate; the case $z=0$ is immediate. Complex
vectors use the real part of the Hermitian pairing in the same argument.
\end{proof}

\KEntry{K03.08.bridge_lebesgue}{Outgoing interface to D11}
\begin{remark}
This is a routing-only record, outside the proved D03 scope. The future
D11 interface must specify its Lebesgue integral and prove that a bounded
Riemann-integrable function on a compact interval is Lebesgue integrable
with the same oriented value. No such assertion is assumed as a lemma
in these notes, and no matching formal-library declaration is claimed.
The countable-cover criterion proved above is an elementary D03 theorem;
it does not silently supply the full Lebesgue integral or this comparison.
\end{remark}
